\documentclass[12pt,aspectratio=1610,professionalfont,serif,t]{beamer}
\usepackage[osf,p]{newpxtext}
\usepackage{newpxmath}
\usepackage{fontawesome5}

\usepackage{fancyvrb}
\fvset{%frame=single,
commandchars=\\\{\},
%framesep=1mm,
fontfamily=helvetica,
fontsize=\normalsize
}
%\usepackage[bitstream-charter]{mathdesign}
\usepackage{listings}
\lstset{commentstyle=\color{orange}, keywordstyle=\color{yellow} }


\usetheme{Boadilla}
\useoutertheme{split}
\usecolortheme{albatross}


\setbeamertemplate{blocks}[rounded][shadow=false]
\setbeamersize{text margin left=\dimexpr 1cm+1ex\relax, text margin right=\dimexpr 1cm+1ex\relax}
\addtobeamertemplate{headline}{}{\usebeamerfont{frametitle}\vspace*{1ex}}
\setbeamertemplate{navigation symbols}{}
\setbeamertemplate{page number in head/foot}{}
%\beamerdefaultoverlayspecification{<+->}


\setbeamercolor*{structure}{fg=green!60!black,bg=black}
\setbeamercolor*{normal text}{fg=green!60!black,bg=black}
\setbeamercolor{palette primary}{use={structure,normal text},fg=green,bg=structure.bg}
\setbeamercolor{palette secondary}{use={structure,normal text},fg=structure.fg,bg=structure.bg}
\setbeamercolor{palette tertiary}{use={structure,normal text},fg=structure.fg,bg=structure.bg}
\setbeamercolor{palette quaternary}{use={structure,normal text},fg=white,bg=green!30!black}
\setbeamercolor{palette sidebar quaternary}{use={structure,normal text},fg=white,bg=green!30!black}
\setbeamercolor{title in head/foot}{use={structure,normal text},fg=white,bg=green!30!black}

\setbeamercolor*{example text}{use={structure},fg=green!60!black}
\setbeamercolor*{block body}{use={structure},bg=structure.bg}
\setbeamercolor*{block body alerted}{use={structure},bg=structure.bg}
\setbeamercolor*{block body example}{use={structure},bg=structure.bg}
\setbeamercolor*{block title}{parent=structure,bg=structure.bg}
\setbeamercolor*{block title alerted}{use={structure,alerted text},fg=alerted text.fg,bg=structure.bg}
\setbeamertemplate{navigation symbols}{}
\setbeamertemplate{items}[square]
\setbeamercolor{item projected}{fg=white}
\setbeamercolor*{normal text}{use={structure},fg=white!90!blue,bg=structure.bg}
\setbeamercolor*{separation line}{}
\setbeamercolor*{fine separation line}{}
\setbeamercolor{alerted text}{use={structure},fg=structure.fg}

\usepackage[italian]{babel}
\usepackage[utf8]{inputenc}
\usepackage{pgf}
\usepackage{verbatim}
\usepackage{inconsolata}
\usepackage{listings}
\lstset{language=C, frame=single,
	basicstyle=\ttfamily,
	numbers=left, numberstyle=\tiny\color{gray},
	numbersep=5pt, fancyvrb=true
}

\usepackage[noend]{algorithm2e}
\usepackage{pgf}
\usepackage{tikz}
\usepackage{graphicx}
\usetikzlibrary{%
	arrows,
	arrows.meta,
	positioning,
	calc,
	backgrounds,
	chains,
	matrix,
	patterns,
	automata,
	fit,
	graphs,
	decorations,
	decorations.pathmorphing,
	decorations.pathreplacing,
	decorations.markings,
}

\usepgflibrary{shapes,shapes.geometric}


\usepackage{url}


\usepackage{xmpmulti}
\usepackage{transparent}
\graphicspath{{img/}}

\newcommand{\zero}{\ensuremath{\mathtt{0}}\xspace}
\newcommand{\uno}{\ensuremath{\mathtt{1}}\xspace}
\newcommand{\due}{\ensuremath{\mathtt{2}}\xspace}
\newcommand{\tre}{\ensuremath{\mathtt{3}}\xspace}
\newcommand{\kleene}{\ensuremath{{}^{\ast}}\xspace}
\xspaceaddexceptions{\zero \uno \due \tre \kleene}


\newcommand{\simplebox}[1]{
      \begin{tikzpicture}
              \node[text centered, rectangle, draw=black, thick, text width=\linewidth,align=center, rounded corners, minimum height=1em] (C) {#1};
      \end{tikzpicture}
}
%
% In questo file si mettono i comandi utili solo per questo insegnamento.
%
\author{Gianluca Della Vedova}
\title{Linguaggi e Computabilità}
\institute{Univ. Milano--Bicocca\\
	\texttt{https://www.unimib.it/gianluca-della-vedova}}
\date{\today}



\begin{document}


\begin{frame}[fragile]
	\frametitle{Linguaggi e Computabilità: Lezione 14}
\begin{enumerate}
\item
Introduzione ai linguaggi liberi dal contesto (CFL) e alle grammatiche libere dal contesto (CFG)
\end{enumerate}
\end{frame}


\begin{frame}
        \frametitle{Esercizio 0}
	Scrivere una grammatica context-free che genera il linguaggio
	\[
                L = \{ \zero^{n} \uno^{n} \; \vert \; n \ge 0 \}
	\]
	sull'alfabeto $\Sigma_{B}$.
% Soluzione: S -> 0 S 1 1 | epsilon
% Ogni applicazione di S -> 0S11 aggiunge uno 0 a sinistra e due 1 a destra.
\end{frame}



\begin{frame}
	\frametitle{Esercizio 1}
	Scrivere una grammatica context-free che genera il linguaggio
	\[
		L = \{ \zero^{n} \uno^{2n} \; \vert \; n \ge 0 \}
	\]
	sull'alfabeto $\Sigma_{B}$.
% Soluzione: S -> 0 S 1 1 | epsilon
% Ogni applicazione di S -> 0S11 aggiunge uno 0 a sinistra e due 1 a destra.
\end{frame}


\begin{frame}
	\frametitle{Esercizio 2}
	Scrivere una grammatica context-free che genera il linguaggio delle
	\emph{palindrome} sull'alfabeto $\Sigma_{B}$, cioè
	\[
		L = \{ w \in \Sigma_{B}\kleene \; \vert \; w = w^{R} \}
	\]
	dove $w^{R}$ è la parola $w$ letta da destra a sinistra.
% Soluzione: S -> 0 S 0 | 1 S 1 | 0 | 1 | epsilon
% Le produzioni S -> 0 | 1 generano le palindrome di lunghezza dispari,
% S -> epsilon quelle di lunghezza pari.
\end{frame}


\begin{frame}
	\frametitle{Esercizio 3}
	Scrivere una grammatica context-free che genera il linguaggio
	\[
		L = \{ \zero^{n} \uno^{m} \due^{m} \tre^{n} \; \vert \; n, m \ge 0 \}
	\]
	sull'alfabeto $\{\zero, \uno, \due, \tre\}$.
% Soluzione:
% S -> 0 S 3 | A
% A -> 1 A 2 | epsilon
% S genera le coppie esterne 0...3, poi A genera le coppie interne 1...2.
\end{frame}


\begin{frame}
	\frametitle{Esercizio 4}
	Scrivere una grammatica context-free che genera il linguaggio delle
	parentesi \emph{bilanciate} sull'alfabeto $\{ (, ) \}$.

	\bigskip
	Esempi: $\varepsilon$, $()$, $(())()$, $((()())())$ appartengono al linguaggio;
	$)($, $(()$, $())($ non appartengono al linguaggio.
% Soluzione: S -> ( S ) S | epsilon
% Ogni parola non vuota bilanciata si scompone in modo unico come (x)y,
% dove ( e ) sono una coppia corrispondente e x, y sono bilanciate.
% Alternativa (ambigua): S -> S S | ( S ) | epsilon
\end{frame}


\begin{frame}
	\frametitle{Esercizio 5}
	Scrivere una grammatica context-free che genera il linguaggio
	\[
		L = \{ w \in \Sigma_{B}\kleene \; \vert \; w \text{ contiene lo stesso numero di \zero e di \uno} \}
	\]
% Soluzione: S -> 0 S 1 S | 1 S 0 S | epsilon
% Una parola non vuota di L inizia con 0 oppure con 1. Se inizia con 0,
% sia 1 il primo simbolo in cui il numero di 0 e di 1 nel prefisso si
% eguaglia: la parola è 0 x 1 y con x, y in L. Analogamente se inizia con 1.
\end{frame}


\end{document}
